\section{Что сломается и по какому признаку это видно}

\begin{table}[htbp]\centering\small
\begin{tabularx}{\textwidth}{@{}p{40mm}Xp{42mm}@{}}
\toprule
Режим отказа & Признак & Реакция \\
\midrule
Совместное правдоподобие не обыгрывает многозадачную базовую линию
& строка 5 абляции внутри интервала парного бутстрапа на нескольких сидах
& зафиксировать как отрицательный результат, отправлять многозадачную модель; решающий и трансдуктивный слои от неё не зависят \\
Голова $E$ как свалка для остатка
& разброс выхода $E$ на соединениях только со скринингом заметно выше, чем на соединениях с кривой
& заменить на константу на фермент, зафиксировать $h_e = 1$ \\
Калибровка зависит от соединения, а не только от фермента
& остатки \eqref{eq:scr} систематически изогнуты по $\pi$ на подмножестве с полными кривыми
& монотонная сплайн-поправка поверх хилловской формы либо смесь двух режимов \\
Серийный приор переобучается на якоря
& на псевдотестовой \emph{страте} (328 соединений, из них 267 с меткой 3A4;
раздел~\ref{sec:proto}) модель с \eqref{eq:series} хуже, чем без. На остальных трёх
ферментах страта слишком мала, и триггер там не срабатывает вовсе
& увеличить $\tau$ или снять слой \\
Конформные интервалы калиброваны маргинально, но не условно
& покрытие по зонам активности расходится с номинальным
& мондриановский конформал с зонами как таксономией \\
Расхождение локальной оценки с промежуточной таблицей
& макро отличается больше чем на 0.03 --- столько даёт переход к тестовой геометрии
на нашей же выборке (0.767 против 0.739 на верхней четверти по сходству)
& чинить разделы \ref{sec:proto} и трансдуктивный слой, а не архитектуру \\
\bottomrule
\end{tabularx}
\caption{Каждому риску --- наблюдаемый признак и заранее решённая реакция.}
\end{table}

Планшетный дрейф был в этом списке и снят измерением (раздел~\ref{sec:calib}).

\section{Что не проверено}

Список на момент выпуска документа.

\emph{Совместная модель запущена в минимальной форме, полная --- нет.} Строка 5 абляции больше
не пуста: общий ствол с двумя головами и переключателем при слагаемом скрининга измерен на четырёх
сидах, результат в разделе~\ref{sec:neg}. Коротко: канал несёт извлекаемую информацию, но
получить её в единицах метрики удаётся только вместе с перекалибровкой после модели, а
поферментно эффект меняет знак и на 2D6 отрицателен.

Незакрытым остаётся следующее. Измерены две свободные головы над общим стволом, то есть
многозадачность; конструкция раздела~\ref{sec:gen}, где скрининг ограничивает $\pi$ через
фиксированную калибровку, --- отдельный объект, и он тоже измерен: результат в
разделе~\ref{sec:neg}, и он хуже, а не лучше.

\emph{Механизм порчи одноголовой руки не установлен, и объясняющая переменная не
идентифицирована.} Шумовая лестница запущена и описана в разделе~\ref{sec:neg}. Она дала
внутрифермент\-ную причинность, которой не было, и сняла информационное объяснение: при почти
мёртвом канале потеря не падает, а растёт. Но чем именно упорядочены ферменты, она различить
не может и не могла: связь скрининга с активностью, ширина доверительной полосы и усиление
калибровочной кривой упорядочивают четыре фермента одинаково, с попарным Спирменом $\pm 1$.
Ширину полосы удалось вычеркнуть пересчётом метрики по тем же предсказаниям с выровненной и с
нулевой полосой; две оставшиеся неразличимы, и шумом их не разделить --- шум двигает
информативность и не двигает ни полосу, ни карту прибора.

Разделяющее вмешательство запущено: калибровочные константы $(E, h)$ обменяны между CYP2D6 и
CYP3A4, всё остальное оставлено прежним. Карта прибора порядок не задаёт --- обменянные
ферменты сдвинулись меньше, чем необменянные контрольные. Остаётся информативность скрининга
либо что-то ещё в данных самого фермента; эксперимент показал, где механизма нет, и не
показал, где он есть. Что могло бы показать, мы пока не придумали.

\emph{Сравнение нейросетевой модели с бустингом пересчитано на четырёх сидах} и больше не
опирается на один испорченный прогон (раздел~\ref{sec:neg}). Итог: в правдоподобном диапазоне
сдвига разница между моделями не превышает 0.005 и ниже $\delta = 0.3$ вообще не измерена,
тогда как постобработка стоит 0.051. Выбор архитектуры здесь почти ничего не решает.

\emph{Само это сравнение оказалось единственным в разделе, сделанным правильно.} Оно считано
после аффинной пары, а абляционная сетка --- до неё, и подраздел~\ref{sec:rawvspost} показал,
чего это стоило: из выигрышей, измеренных сырыми предсказаниями, после пары выживают два ---
канал скрининга в стволе и механистический блок, ровно те, у которых прирост ранговой связи с
истиной отличен от нуля. Причём блок сырое измерение не завышало, а занижало, и он два месяца
числился отрицательным результатом. Проверять новые руки теперь надо в той же точке, где считается результат, и
печатать рядом прирост ранга: он и говорит заранее, переживёт ли рука постобработку.

Предобученный энкодер запущен, и в двух видах: как замена нашему представлению и как добавка
к нему (раздел~\ref{sec:neg}). Не запускались коррегионализация и совместное правдоподобие
поверх неё, то есть строка 2 и полная версия 5. Это по-прежнему объясняет, почему конструкция
должна работать, а не показывает, что она работает в полном виде.

\emph{Вероятностное правило TDI} --- интегрирование \eqref{eq:tdirule} по совместному предсказательному
распределению --- не запускалось; все измеренные варианты детерминированные.

\emph{Оценщик $\mathrm{p}K_a$} проверен на двадцати соединениях, выбранных вручную. Средняя ошибка 0.60
при медианной 0.2, то есть среднее делают четыре промаха, худший --- кофеин, 4.5 единицы: правило относит
его азот к имидазолам с базовым 7.0 и недоштрафовывает соседние карбонилы. На бинарный вопрос про заряд
при 7.4 это не влияет, но мотив «ароматический азот через атом от карбонила» встречается у трёх процентов
обучающей выборки --- зона, где правило может систематически завышать.

Проверка, которая это меряет, до сих пор не запускалась: у неё было сломано две подстановки путей и не
хватало каталога в \texttt{sys.path}, а внутри она читала не то поле возврата --- индекс атома вместо
имени класса, из-за чего два её числа печатались нулями. После починки картина шире заявленной.
Самый основной центр относится к имидазолу или пиридину у 55.8 процента обучающей выборки, то есть
больше чем у половины молекул ответ правила опирается на ароматический азот. Пересечение с карбонилом
по соседству --- ровно конфигурация кофеина --- даёт 109 молекул, 2.2 процента. Три процента из
предыдущего абзаца остаются верными, но относятся к мотиву, а не к тем молекулам, у которых на этом
мотиве стоит итоговый ответ.

\emph{Учит ли графовая сеть протонирование сама} --- не проверялось. Если учит, часть выводов
про механистический блок придётся переписать.

\emph{Зависимость калибровки от соединения} в этих данных непроверяема: у каждого соединения одно показание
скрининга, разделить «соединение отклоняется от калибровки» и «прибор шумит» на одном измерении нельзя.

\emph{Контроль в проверке ограничения диапазона} выравнивает количество примеров, но не форму
распределения меток.

\emph{Механизм пулирования не установлен, и это худший вид незакрытости.} Пулирование --- самый
крупный эффект, на который опирается сабмит, и оба естественных объяснения проверены и
опровергнуты.

Первое: пулированная модель дотягивается до соседей, размеченных по \emph{другому} ферменту,
которых поферментная модель не видит вовсе. Матрица меток разрежена, невидимая доля окрестности
велика --- от 20 до 45 процентов строк имеют зазор больше 0.02 между ближайшим соседом вообще и
ближайшим размеченным по этому ферменту. Но преимущество пула по величине зазора идёт не туда:
$+0.0165$ там, где зазор нулевой и одалживать нечего, и немонотонно дальше. Механизм требует нуля
в первой ячейке; он его не даёт.

Второе: пул переносит общую часть зависимости структура--активность. Тогда выигрывать должен
фермент, наиболее связанный с остальными. Связь оказалась перевёрнутой без единого исключения,
$\rho = -1.000$ по четырём точкам: CYP2D6 со средней корреляцией $-0.002$ берёт $+0.0374$, а
CYP3A4 со средней $0.375$ --- единственный, кому пул вредит ($-0.0057$). Пул выигрывает ровно
там, где общего нет.

Остаётся третье прочтение --- не сходство, а \emph{контраст}: индикатор фермента позволяет дереву
выучить <<этот сплит важен для 2D6 и не важен для 3A4>>, чего поферментная модель не может
выразить в принципе, и тогда полезность партнёра растёт с тем, насколько он \emph{отличается}.
Оно согласуется со знаком, но не проверено, и отличить его от простой регуляризации числом строк
по четырём точкам нельзя: у CYP3A4 и меток больше всех, и корреляция выше всех.

Формулировать надо жёстко. Необъяснённый эффект такого размера опаснее отрицательного: мы не
знаем, от чего он зависит, и потому не можем утверждать, что он перенесётся на тест. Разделяющие
прогоны --- попарное пулирование и пул с обнулённым индикатором --- поставлены и на момент выпуска
не досчитаны.

\section{Сводка чисел}

\begin{table}[htbp]\centering\small
\begin{tabularx}{\textwidth}{@{}Xrrrrr@{}}
\toprule
& 1A2 & 2C9 & 2D6 & 3A4 & Макро \\
\midrule
Меток $\pic$ & 1412 & 1285 & 1493 & 2335 & 6525 \\
Межквартильный размах меток & 1.00 & 0.92 & 0.83 & 1.50 & --- \\
Надёжность меток (потолок $R^2$) & 0.948 & 0.905 & 0.934 & 0.898 & --- \\
\addlinespace[2pt]
ST-RAE, фингерпринт и дескрипторы & 0.880 & 0.702 & 0.993 & 0.518 & 0.773 \\
ST-RAE, плюс механистический блок & 0.879 & 0.691 & 0.980 & 0.519 & 0.767 \\
Спирмен, фингерпринт и дескрипторы & 0.497 & 0.583 & 0.351 & 0.764 & 0.549 \\
Спирмен, плюс механистический блок & 0.496 & 0.597 & 0.403 & 0.765 & 0.565 \\
\addlinespace[2pt]
Калибровка: $E$ & 0.728 & 0.621 & 0.867 & 0.931 & --- \\
Калибровка: $h$ & 1.261 & 1.112 & 1.243 & 1.968 & --- \\
Ранговая связь скрининга с $\pic$ & $-0.86$ & $-0.90$ & $-0.83$ & $-0.94$ & --- \\
\addlinespace[2pt]
Ско ST-RAE по четырём сидам, без блока & 0.0136 & 0.0098 & 0.0132 & 0.0063 & 0.0066 \\
Ско ST-RAE по четырём сидам, с блоком & 0.0109 & 0.0078 & 0.0052 & 0.0057 & 0.0019 \\
\bottomrule
\end{tabularx}
\caption{Кросс-валидированы шесть строк из двенадцати. Четыре строки ST-RAE и Спирмена --- пятикратная
кросс-валидация по кластерам Бутины при пороге 0.35, макро на сиде 0; две строки со стандартным
отклонением --- по четырём сидам того же разбиения. Остальные шесть кросс-валидацией не являются и
сравнивать их с первыми по строгости нельзя: число меток и межквартильный размах --- статистики колонки,
надёжность --- разложение дисперсии по заявленной погрешности измерения, а калибровка $E$, $h$ и ранговая
связь скрининга --- подгонки на всей выборке сразу, без отложенной части.}
\end{table}

\begin{table}[htbp]\centering\small
\begin{tabularx}{0.82\textwidth}{@{}Xrr@{}}
\toprule
Правило решения для TDI & MCC 3A4 & MCC 2D6 \\
\midrule
Всегда «нет» & 0.000 & 0.000 \\
Классификатор, порог 0.5 & 0.302 & 0.103 \\
Классификатор, порог подобран & 0.356 & 0.127 \\
Две независимые регрессии плеч, затем правило & 0.255 & 0.094 \\
Отдельная голова на сдвиг, затем правило & 0.323 & 0.104 \\
То же с подобранным порогом сдвига & 0.324 & 0.149 \\
\addlinespace[2pt]
Подстановочный порог по сырым вероятностям & 0.325 & 0.115 \\
Подстановочный порог после калибровки Платта & 0.352 & 0.092 \\
\bottomrule
\end{tabularx}
\caption{Все числа получены на 2346 соединениях для 3A4 и 1495 для 2D6 --- тех, у кого есть и признаки,
и метка. Доля положительных в этих выборках 32.6 и 21.7 процента; по файлу TDI целиком --- 21.3 и 21.6,
и разница целиком объясняется тем, что 1238 соединений без кривой доза-эффект автоматически размечены
как отрицательные.}
\end{table}
